
\documentclass[12pt,a4paper]{article}

\usepackage[T1]{fontenc}
\usepackage{lmodern}
\usepackage{graphicx}
\usepackage{array}

\title{Tables and Figures in LaTeX}
\author{Vikram Mehra}
\date{\today}

\begin{document}

\maketitle

\section{Introduction}

LaTeX provides environments for creating professional
tables and inserting figures into documents.

The tabular environment is used to arrange information
in rows and columns. The graphicx package is used to
insert external images.

\section{Creating a Basic Table}

The following table demonstrates column alignment
and horizontal borders.

\begin{center}
\begin{tabular}{|c|l|c|}
\hline
Roll Number & Student Name & Marks \\
\hline
101 & Student A & 85 \\
102 & Student B & 90 \\
103 & Student C & 88 \\
\hline
\end{tabular}
\end{center}

\section{Table with Caption and Label}

Tables can be placed inside the table environment.
The caption describes the table, while the label
allows us to refer to it elsewhere.

\begin{table}[htbp]
\centering
\begin{tabular}{|c|l|c|}
\hline
Subject & Description & Marks \\
\hline
Mathematics & Algebra and Calculus & 90 \\
Physics & Mechanics and Optics & 85 \\
Computer Science & Programming & 95 \\
\hline
\end{tabular}
\caption{Example of student subject marks}
\label{tab:marks}
\end{table}

As shown in Table~\ref{tab:marks}, LaTeX can
automatically number tables and create references.

\section{Combining Table Columns}

The multicolumn command combines multiple columns.

\begin{table}[htbp]
\centering
\begin{tabular}{|c|c|c|}
\hline
\multicolumn{3}{|c|}{Student Information} \\
\hline
Roll Number & Name & Grade \\
\hline
101 & Student A & A \\
102 & Student B & B \\
\hline
\end{tabular}
\caption{Table using the multicolumn command}
\label{tab:multicolumn}
\end{table}

\section{Inserting a Figure}

The figure environment is used to insert and manage
images in a document.

The following box is a placeholder for an image.

\begin{figure}[htbp]
\centering
\fbox{\rule{0pt}{4cm}\rule{0.75\linewidth}{0pt}}
\caption{Placeholder for an example figure}
\label{fig:placeholder}
\end{figure}

Figure~\ref{fig:placeholder} demonstrates how to
add a caption and reference to a figure.

\section{Using the Includegraphics Command}

To insert an actual image, first place an image file
named sample-image.png in the same folder as this
source file.

Then replace the placeholder figure above with:

% \begin{figure}[htbp]
% \centering
% \includegraphics[width=0.7\textwidth]{sample-image.png}
% \caption{Example of an inserted image}
% \label{fig:sample}
% \end{figure}

\section{Conclusion}

This program demonstrates tables, column alignment,
captions, labels, references, and figure insertion
using LaTeX.

\end{document}
